\documentclass[aspectratio=169]{beamer}

% =========================================================
% TEMPLATE INSTITUCIONAL
% =========================================================
% Compilar con XeLaTeX/LuaLaTeX y -shell-escape si se usan bloques minted.
% =========================================================
% template.tex  (NO tiene \documentclass ni \begin{document})
    % Plantilla institucional robusta: modo claro,
    % fallbacks de imágenes, minted, entornos de código claros/oscuros y macros docentes.
    % =========================================================

    % -------------------------
    % CONFIG POR DEFECTO
    % (En main.tex puedes sobrescribir con \renewcommand)
    % -------------------------
    \newcommand{\CourseName}{Laboratorio Inteligencia Artificial}
    \newcommand{\CourseCode}{ING01336}
    \newcommand{\Faculty}{Facultad de Ingenierías}
    \newcommand{\Institution}{Politécnico Colombiano Jaime Isaza Cadavid}
    \newcommand{\Instructor}{Deimer Miranda Montoya, MSc.(c).}
    \newcommand{\Email}{deimer\_miranda91162@elpoli.edu.co}
    \newcommand{\InstitutionLogo}{Template/LogoSena.png}

    % Control global de imágenes: 1 = mostrar, 0 = no cargar
    \newcommand{\UseImages}{1}

    % -------------------------
    % MOTOR Y FUENTES (ROBUSTO)
    % -------------------------
    \usepackage{iftex}
    \ifPDFTeX
    \usepackage[utf8]{inputenc}
    \usepackage[T1]{fontenc}
    \usepackage{lmodern}
    \else
    \usepackage{fontspec}
    \usepackage[sfdefault]{noto}
    \fi



    % -------------------------
    % PAQUETES
    % -------------------------
    \usetheme{Madrid}
    \usepackage{xparse}
    \usepackage{fontawesome5}
    \usepackage{booktabs}
    \usepackage{xcolor}
    \usepackage{graphicx}
    \usepackage{tikz}
    \usetikzlibrary{calc,positioning,arrows.meta,mindmap,shapes.geometric}
    \usepackage{hyperref}
    \usepackage{tcolorbox}
    \tcbuselibrary{minted,skins,breakable}

    \hypersetup{colorlinks=true, linkcolor=blue, urlcolor=blue, citecolor=blue}

    % -------------------------
    % COLORES (INSTITUCIONAL)
    % -------------------------
    \definecolor{BrandYellow}{RGB}{255,185,62}
    \definecolor{BrandGreen}{RGB}{20,196,134}
    \definecolor{BrandBlue}{RGB}{30,60,120}
    \definecolor{DarkBG}{RGB}{14,16,20}
    \definecolor{DarkFG}{RGB}{235,235,235}

    \definecolor{SoftGray}{RGB}{245,246,248}
    \definecolor{MediumGray}{RGB}{210,213,218}
    \definecolor{DarkGray}{RGB}{45,47,52}

    % Fondos de código
    \definecolor{pybg}{rgb}{0.12,0.18,0.30}
    \definecolor{cppbg}{rgb}{0.10,0.10,0.10}
    \definecolor{bashbg}{rgb}{0.2,0.2,0.2}

    % Fondos específicos para código HTML
    \definecolor{HtmlClaroFondo}{HTML}{F7F7F7}
    \definecolor{HtmlClaroBorde}{HTML}{D0D0D0}
    \definecolor{HtmlOscuroFondo}{HTML}{1E1E1E}
    \definecolor{HtmlOscuroBorde}{HTML}{444444}

    % -------------------------
    % COLORES PARA BLOQUES (DIFERENTES)
    % -------------------------
    % Idea
    \definecolor{IdeaTitle}{RGB}{30,60,120}     % azul
    \definecolor{IdeaBody}{RGB}{235,241,255}    % azul muy claro

    % Practice
    \definecolor{PracticeTitle}{RGB}{20,120,90} % verde/teal
    \definecolor{PracticeBody}{RGB}{232,250,245}

    % Definition
    \definecolor{DefTitle}{RGB}{90,70,160}      % violeta
    \definecolor{DefBody}{RGB}{240,236,255}

    % Result
    \definecolor{ResultTitle}{RGB}{0,120,0}     % verde
    \definecolor{ResultBody}{RGB}{235,250,235}

    % Warning (alert)
    \definecolor{WarnTitle}{RGB}{160,30,30}     % rojo
    \definecolor{WarnBody}{RGB}{255,235,235}

    % Comparación
    \definecolor{CompareTitle}{RGB}{28,92,145}
    \definecolor{CompareBody}{RGB}{235,245,252}

    % Pregunta
    \definecolor{QuestionTitle}{RGB}{176,112,20}
    \definecolor{QuestionBody}{RGB}{255,247,229}

    % Ejemplo
    \definecolor{ExampleTitle}{RGB}{84,78,150}
    \definecolor{ExampleBody}{RGB}{243,241,255}

    % =========================================================
    % TEMA VISUAL - SOLO MODO CLARO
    % =========================================================
    \setbeamercolor{background canvas}{bg=white}
    \setbeamercolor{normal text}{fg=black,bg=white}
    \setbeamercolor{structure}{fg=BrandBlue}
    \setbeamercolor{itemize item}{fg=BrandBlue}
    \setbeamercolor{itemize subitem}{fg=BrandBlue}
    \setbeamercolor{enumerate item}{fg=BrandBlue}
    \setbeamercolor{alerted text}{fg=WarnTitle}
    \setbeamercolor{example text}{fg=PracticeTitle}
    \setbeamercolor{caption name}{fg=BrandBlue}

    % -------------------------
    % Fondo tipo banda superior (frametitle)
    % -------------------------
    \setbeamercolor{frametitle}{fg=white,bg=BrandBlue}
    \setbeamertemplate{frametitle}{
        \nointerlineskip
        \begin{beamercolorbox}[wd=\paperwidth,ht=0.5cm,dp=0.2cm,leftskip=0.6cm]{frametitle}
            \usebeamerfont{frametitle}\insertframetitle
        \end{beamercolorbox}
    }

    % -------------------------
    % IMÁGENES (ROBUSTAS)
    % -------------------------
    \graphicspath{{Template/}{Images/}{Images2/}{Logos/}}

    \newcommand{\smartimage}[2][0.9\linewidth]{%
        \ifnum\UseImages=1
        \IfFileExists{#2}{%
            \includegraphics[width=#1]{#2}%
        }{%
            \begingroup
            \setlength{\fboxsep}{4pt}%
            \fbox{\parbox{#1}{\centering\scriptsize Imagen no encontrada:\\ \texttt{#2}}}%
            \endgroup
        }%
        \fi
    }

    \newcommand{\smartlogo}[2][1.8cm]{%
        \ifnum\UseImages=1
        \IfFileExists{#2}{%
            \includegraphics[height=#1]{#2}%
        }{%
            {\footnotesize\textbf{[Logo no encontrado]}}%
        }%
        \fi
    }

    % -------------------------
    % CÓDIGO (MINTED)
    % Se conservan entornos claros y oscuros de código independientemente del tema visual.
    % Requiere compilar con -shell-escape y pygments instalado
    % -------------------------
    \usepackage{minted}
    \usemintedstyle{monokai}

    % =========================================================
    % COLORES PARA ENTORNOS DE CÓDIGO
    % =========================================================

    % Fondo claro general
    \definecolor{CodigoClaroFondo}{HTML}{F7F7F7}
    \definecolor{CodigoClaroBorde}{HTML}{D0D0D0}

    % Fondo oscuro general
    \definecolor{CodigoOscuroFondo}{HTML}{1E1E1E}
    \definecolor{CodigoOscuroBorde}{HTML}{444444}

    % =========================================================
    % ENTORNOS HTML CON RESALTADO DE SINTAXIS
    % Todos los frames que los utilicen deben llevar [fragile].
    % =========================================================

    % Fragmentos cortos HTML con fondo claro y sin numeración
    \newtcblisting{htmlcode}{
        listing engine=minted,
        minted language=html,
        minted style=friendly,
        minted options={
            fontsize=\small,
            breaklines=true,
            autogobble=true,
            stripnl=true,
            tabsize=4,
            encoding=utf8
        },
        colback=HtmlClaroFondo,
        colframe=HtmlClaroBorde,
        coltext=black,
        boxrule=0.6pt,
        arc=2mm,
        left=2mm,
        right=2mm,
        top=0.0mm,
        bottom=0.0mm,
        before skip=1pt,
        after skip=1pt,
        listing only,
        enhanced
    }

    % Fragmentos cortos HTML con fondo oscuro y sin numeración
    \newtcblisting{htmlcodedark}{
        listing engine=minted,
        minted language=html,
        minted style=monokai,
        minted options={
            fontsize=\small,
            breaklines=true,
            autogobble=true,
            stripnl=true,
            tabsize=4,
            encoding=utf8
        },
        colback=HtmlOscuroFondo,
        colframe=HtmlOscuroBorde,
        coltext=white,
        boxrule=0.6pt,
        arc=2mm,
        left=2mm,
        right=2mm,
        top=0.0mm,
        bottom=0.0mm,
        before skip=1pt,
        after skip=1pt,
        listing only,
        enhanced
    }

    % =========================================================
    % PYTHON
    % =========================================================

    % Python con fondo claro
    \newtcblisting{pythoncode}{
        listing engine=minted,
        minted language=python,
        minted style=friendly,
        minted options={
            fontsize=\small,
            breaklines=true,
            autogobble=true,
            stripnl=true,
            tabsize=4,
            encoding=utf8
        },
        colback=CodigoClaroFondo,
        colframe=CodigoClaroBorde,
        coltext=black,
        boxrule=0.6pt,
        arc=2mm,
        left=2mm,
        right=2mm,
        top=0.0mm,
        bottom=0.0mm,
        before skip=1pt,
        after skip=1pt,
        listing only,
        enhanced
    }

    % Python con fondo oscuro
    \newtcblisting{pythoncodedark}{
        listing engine=minted,
        minted language=python,
        minted style=monokai,
        minted options={
            fontsize=\small,
            breaklines=true,
            autogobble=true,
            stripnl=true,
            tabsize=4,
            encoding=utf8
        },
        colback=CodigoOscuroFondo,
        colframe=CodigoOscuroBorde,
        coltext=white,
        boxrule=0.6pt,
        arc=2mm,
        left=2mm,
        right=2mm,
        top=0.0mm,
        bottom=0.0mm,
        before skip=1pt,
        after skip=1pt,
        listing only,
        enhanced
    }

    % =========================================================
    % C++
    % =========================================================

    % C++ con fondo claro
    \newtcblisting{cppcode}{
        listing engine=minted,
        minted language=cpp,
        minted style=friendly,
        minted options={
            fontsize=\small,
            breaklines=true,
            autogobble=true,
            stripnl=true,
            tabsize=4,
            encoding=utf8
        },
        colback=CodigoClaroFondo,
        colframe=CodigoClaroBorde,
        coltext=black,
        boxrule=0.6pt,
        arc=2mm,
        left=2mm,
        right=2mm,
        top=0.0mm,
        bottom=0.0mm,
        before skip=1pt,
        after skip=1pt,
        listing only,
        enhanced
    }

    % C++ con fondo oscuro
    \newtcblisting{cppcodedark}{
        listing engine=minted,
        minted language=cpp,
        minted style=monokai,
        minted options={
            fontsize=\small,
            breaklines=true,
            autogobble=true,
            stripnl=true,
            tabsize=4,
            encoding=utf8
        },
        colback=CodigoOscuroFondo,
        colframe=CodigoOscuroBorde,
        coltext=white,
        boxrule=0.6pt,
        arc=2mm,
        left=2mm,
        right=2mm,
        top=0.0mm,
        bottom=0.0mm,
        before skip=1pt,
        after skip=1pt,
        listing only,
        enhanced
    }

    % =========================================================
    % CSS
    % =========================================================

    % CSS con fondo claro
    \newtcblisting{csscode}{
        listing engine=minted,
        minted language=css,
        minted style=friendly,
        minted options={
            fontsize=\small,
            breaklines=true,
            autogobble=true,
            stripnl=true,
            tabsize=4,
            encoding=utf8
        },
        colback=CodigoClaroFondo,
        colframe=CodigoClaroBorde,
        coltext=black,
        boxrule=0.6pt,
        arc=2mm,
        left=2mm,
        right=2mm,
        top=0.0mm,
        bottom=0.0mm,
        listing only,
        enhanced
    }

    % CSS con fondo oscuro
    \newtcblisting{csscodedark}{
        listing engine=minted,
        minted language=css,
        minted style=monokai,
        minted options={
            fontsize=\small,
            breaklines=true,
            autogobble=true,
            stripnl=true,
            tabsize=4,
            encoding=utf8
        },
        colback=CodigoOscuroFondo,
        colframe=CodigoOscuroBorde,
        coltext=white,
        boxrule=0.6pt,
        arc=2mm,
        left=2mm,
        right=2mm,
        top=0.0mm,
        bottom=0.0mm,
        listing only,
        enhanced
    }

    % =========================================================
    % SQL
    % =========================================================

    % SQL con fondo claro
    \newtcblisting{sqlcode}{
        listing engine=minted,
        minted language=sql,
        minted style=friendly,
        minted options={
            fontsize=\small,
            breaklines=true,
            autogobble=true,
            stripnl=true,
            tabsize=2,
            encoding=utf8
        },
        colback=CodigoClaroFondo,
        colframe=CodigoClaroBorde,
        coltext=black,
        boxrule=0.6pt,
        arc=2mm,
        left=2mm,
        right=2mm,
        top=0.0mm,
%        before skip=1pt,
%        after skip=1pt,
        bottom=0.0mm,
        listing only,
        enhanced
    }

    % SQL con fondo oscuro
    \newtcblisting{sqlcodedark}{
        listing engine=minted,
        minted language=sql,
        minted style=monokai,
        minted options={
            fontsize=\small,
            breaklines=true,
            autogobble=true,
            stripnl=true,
            tabsize=2,
            encoding=utf8
        },
        colback=CodigoOscuroFondo,
        colframe=CodigoOscuroBorde,
        coltext=white,
        boxrule=0.6pt,
        arc=2mm,
        left=2mm,
        right=2mm,
        top=0.0mm,
        bottom=0.0mm,
%        before skip=1pt,
%        after skip=1pt,
        listing only,
        enhanced
    }

    % =========================================================
    % JAVA
    % =========================================================

    % Java con fondo claro
    \newtcblisting{javacode}{
        listing engine=minted,
        minted language=java,
        minted style=friendly,
        minted options={
            fontsize=\small,
            breaklines=true,
            autogobble=true,
            stripnl=true,
            tabsize=4,
            encoding=utf8
        },
        colback=CodigoClaroFondo,
        colframe=CodigoClaroBorde,
        coltext=black,
        boxrule=0.6pt,
        arc=2mm,
        left=2mm,
        right=2mm,
        top=0.0mm,
        bottom=0.0mm,
%        before skip=1pt,
%        after skip=1pt,
        listing only,
        enhanced
    }

    % Java con fondo oscuro
    \newtcblisting{javacodedark}{
        listing engine=minted,
        minted language=java,
        minted style=monokai,
        minted options={
            fontsize=\small,
            breaklines=true,
            autogobble=true,
            stripnl=true,
            tabsize=4,
            encoding=utf8
        },
        colback=CodigoOscuroFondo,
        colframe=CodigoOscuroBorde,
        coltext=white,
        boxrule=0.6pt,
        arc=2mm,
        left=2mm,
        right=2mm,
        top=0.0mm,
        bottom=0.0mm,
%        before skip=1pt,
%        after skip=1pt,
        listing only,
        enhanced
    }

    % =========================================================
    % BASH
    % =========================================================

    % Bash con fondo claro
    \newtcblisting{bashcode}{
        listing engine=minted,
        minted language=bash,
        minted style=friendly,
        minted options={
            fontsize=\small,
            breaklines=true,
            autogobble=true,
            stripnl=true,
            tabsize=4,
            encoding=utf8
        },
        colback=CodigoClaroFondo,
        colframe=CodigoClaroBorde,
        coltext=black,
        boxrule=0.6pt,
        arc=2mm,
        left=2mm,
        right=2mm,
        top=0.0mm,
        bottom=0.0mm,
%        before skip=1pt,
%        after skip=1pt,
        listing only,
        enhanced
    }

    % Bash con fondo oscuro
    \newtcblisting{bashcodedark}{
        listing engine=minted,
        minted language=bash,
        minted style=monokai,
        minted options={
            fontsize=\small,
            breaklines=true,
            autogobble=true,
            stripnl=true,
            tabsize=4,
            encoding=utf8
        },
        colback=CodigoOscuroFondo,
        colframe=CodigoOscuroBorde,
        coltext=white,
        boxrule=0.6pt,
        arc=2mm,
        left=2mm,
        right=2mm,
        top=0.0mm,
        bottom=0.0mm,
%        before skip=1pt,
%        after skip=1pt,
        listing only,
        enhanced
    }

    % =========================================================
    % MONGODB / MONGOSH - FONDO CLARO
    % =========================================================

    \newtcblisting{mongodbcode}{
        listing engine=minted,
        minted language=javascript,
        minted style=friendly,
        minted options={
            fontsize=\small,
            breaklines=true,
            autogobble=true,
            stripnl=true,
            tabsize=4,
            encoding=utf8
        },
        colback=CodigoClaroFondo,
        colframe=CodigoClaroBorde,
        coltext=black,
        boxrule=0.6pt,
        arc=2mm,
        left=2mm,
        right=2mm,
        top=0.0mm,
        bottom=0.0mm,
        listing only,
        enhanced,
        breakable
    }

    % =========================================================
    % MONGODB / MONGOSH - FONDO OSCURO
    % =========================================================

    \newtcblisting{mongodbcodedark}{
        listing engine=minted,
        minted language=javascript,
        minted style=monokai,
        minted options={
            fontsize=\small,
            breaklines=true,
            autogobble=true,
            stripnl=true,
            tabsize=4,
            encoding=utf8
        },
        colback=CodigoOscuroFondo,
        colframe=CodigoOscuroBorde,
        coltext=white,
        boxrule=0.6pt,
        arc=2mm,
        left=2mm,
        right=2mm,
        top=0.0mm,
        bottom=0.0mm,
        listing only,
        enhanced,
        breakable
    }

    % =========================================================
    % CONSOLA / SALIDA DE TERMINAL - TEMA OSCURO
    % =========================================================
    \newtcblisting{consolecodedark}{
        listing engine=minted,
        minted language=text,
        minted options={
            breaklines,
            stripnl=true,
            breakanywhere,
            fontsize=\small,
            autogobble
        },
        listing only,
        enhanced,
        breakable,
        colback=black!92,
        colframe=black!65,
        coltext=white,
        boxrule=0.6pt,
        arc=2mm,
        left=2mm,
        right=2mm,
        top=0.0mm,
        bottom=0.0mm,
        before skip=2mm,
        after skip=2mm
    }
    % =========================================================
    % CONSOLA / SALIDA DE TERMINAL - TEMA CLARO
    % =========================================================
    \newtcblisting{consolecode}{
        listing engine=minted,
        minted language=text,
        minted options={
            breaklines,
            stripnl=true,
            breakanywhere,
            fontsize=\small,
            autogobble
        },
        listing only,
        enhanced,
        breakable,
        colback=black!4,
        colframe=black!30,
        coltext=black,
        boxrule=0.6pt,
        arc=2mm,
        left=2mm,
        right=2mm,
        top=0.0mm,
        bottom=0.0mm,
        before skip=2mm,
        after skip=2mm
    }

    % -------------------------
    % MACROS DOCENTES
    % -------------------------
    \newcommand{\SectionSlide}[2]{%
        \begin{frame}
            \vspace*{0.1cm}
            \centering
            \begin{tikzpicture}
                \node[fill=BrandBlue, rounded corners,
                minimum width=0.86\textwidth, minimum height=1.6cm,
                text width=0.82\textwidth, align=center] (box)
                {\color{white}\textbf{\huge #1}\\[2pt]\Large #2};
            \end{tikzpicture}
        \end{frame}
    }

    % =========================================================
    % BLOQUES DOCENTES CON TÍTULO OPCIONAL + COLORES DIFERENTES
    % =========================================================
    % Idea: \Idea[Mi título]{Contenido}
    \NewDocumentCommand{\Idea}{O{Idea clave} +m}{
        \begingroup
        \setbeamercolor{block title}{fg=white,bg=IdeaTitle}
        \setbeamercolor{block body}{fg=black,bg=IdeaBody}
        \begin{block}{\faLightbulb\ \,#1}
            #2
        \end{block}
        \endgroup
    }

    % Practice
    \NewDocumentCommand{\Practice}{O{Actividad práctica} +m}{
        \begingroup
        \setbeamercolor{block title}{fg=white,bg=PracticeTitle}
        \setbeamercolor{block body}{fg=black,bg=PracticeBody}
        \begin{block}{\faLaptopCode\ \,#1}
            #2
        \end{block}
        \endgroup
    }

    % Definition
    \NewDocumentCommand{\DefBlock}{O{Definición} +m}{
        \begingroup
        \setbeamercolor{block title}{fg=white,bg=DefTitle}
        \setbeamercolor{block body}{fg=black,bg=DefBody}
        \begin{block}{\faBook\ \,#1}
            #2
        \end{block}
        \endgroup
    }

    % Result
    \NewDocumentCommand{\Result}{O{Resultado} +m}{
        \begingroup
        \setbeamercolor{block title}{fg=white,bg=ResultTitle}
        \setbeamercolor{block body}{fg=black,bg=ResultBody}
        \begin{block}{\faCheckCircle\ \,#1}
            #2
        \end{block}
        \endgroup
    }

    % Warning (alertblock)
    \NewDocumentCommand{\Warning}{O{Atención} +m}{
        \begingroup
        \setbeamercolor{block title alerted}{fg=white,bg=WarnTitle}
        \setbeamercolor{block body alerted}{fg=black,bg=WarnBody}
        \begin{alertblock}{\faExclamationTriangle\ \,#1}
            #2
        \end{alertblock}
        \endgroup
    }

    % Comparación: \Compare[Mi título]{Contenido}
    \NewDocumentCommand{\Compare}{O{Comparación} +m}{
        \begingroup
        \setbeamercolor{block title}{fg=white,bg=CompareTitle}
        \setbeamercolor{block body}{fg=black,bg=CompareBody}
        \begin{block}{\faBalanceScale\ \,#1}
            #2
        \end{block}
        \endgroup
    }

    % Pregunta: \Question[Mi título]{Contenido}
    \NewDocumentCommand{\Question}{O{Pregunta guía} +m}{
        \begingroup
        \setbeamercolor{block title}{fg=white,bg=QuestionTitle}
        \setbeamercolor{block body}{fg=black,bg=QuestionBody}
        \begin{block}{\faQuestionCircle\ \,#1}
            #2
        \end{block}
        \endgroup
    }

    % Ejemplo: \ExampleBlock[Mi título]{Contenido}
    % Se usa ExampleBlock para evitar colisión con comandos propios de Beamer.
    \NewDocumentCommand{\ExampleBlock}{O{Ejemplo} +m}{
        \begingroup
        \setbeamercolor{block title}{fg=white,bg=ExampleTitle}
        \setbeamercolor{block body}{fg=black,bg=ExampleBody}
        \begin{block}{\faPuzzlePiece\ \,#1}
            #2
        \end{block}
        \endgroup
    }

    \newcommand{\AgendaSlide}{%
        \begin{frame}{Agenda}
            \tableofcontents
        \end{frame}
    }

    % =========================================================
    % SLIDE DE NORMAS INSTITUCIONALES
    % =========================================================
    \newcommand{\NormasSlide}{%
        \begin{frame}{Normas del curso}
            \begin{minipage}{0.32\textwidth}
                \centering
                \smartimage[1cm]{Template/Puntual.png}\\[0.2cm]
                \textbf{Cumple con los horarios y demuestra compromiso académico en cada sesión.}
            \end{minipage}\hfill
            \begin{minipage}{0.32\textwidth}
                \centering
                \smartimage[1cm]{Template/Actitud.png}\\[0.2cm]
                \textbf{Mantén una actitud profesional, ética y orientada al aprendizaje continuo.}
            \end{minipage}\hfill
            \begin{minipage}{0.32\textwidth}
                \centering
                \smartimage[1cm]{Template/Comunicacion.png}\\[0.2cm]
                \textbf{Comunica tus ideas con claridad, respeto y fundamento técnico.}
            \end{minipage}

            \vspace{0.2cm}

            \begin{minipage}{0.48\textwidth}
                \centering
                \smartimage[1cm]{Template/Participacion.png}\\[0.2cm]
                \textbf{Participa activamente y contribuye al trabajo colaborativo.}
            \end{minipage}\hfill
            \begin{minipage}{0.48\textwidth}
                \centering
                \smartimage[1cm]{Template/Cuidado.png}\\[0.2cm]
                \textbf{Haz uso responsable de los recursos, equipos e instalaciones.}
            \end{minipage}
        \end{frame}
    }


    % -------------------------
    % PIE DE PÁGINA
    % -------------------------
    \setbeamertemplate{navigation symbols}{}

    \setbeamertemplate{footline}{
        \leavevmode
        \hbox to \paperwidth{%
            \hspace{0.35cm}%
            % Zona izquierda: navegación y numeración
            \makebox[0.22\paperwidth][l]{%
                {\tiny
                    \Acrobatmenu{PrevPage}{\textcolor{black}{$\triangleleft$}}%
                    \hspace{0.22cm}%
                    \insertpagenumber/\insertpresentationendpage%
                    \hspace{0.22cm}%
                    \Acrobatmenu{NextPage}{\textcolor{black}{$\triangleright$}}%
                }%
            }%
            \hfill%
            % Zona derecha: instructor + logo
            \makebox[0.70\paperwidth][r]{%
                {\tiny\Instructor}%
                \hspace{0.18cm}%
                \raisebox{-0.10cm}{\smartlogo[0.34cm]{\InstitutionLogo}}%
                %            \raisebox{-0.10cm}{\smartlogo[0.34cm]{Template/LogoSena.png}}%
            }%
            \hspace{0.35cm}%
        }%
        \vspace{0.08cm}
    }
    % -------------------------
    % PORTADA (UN COMANDO)
    % -------------------------
    \newcommand{\TitleSlide}{%
        \begin{frame}
            \centering
            %\smartlogo[1.8cm]{logo.png}

            \vspace{0.4cm}
            \color{BrandGreen}\rule{0.9\textwidth}{0.5pt}

            \vspace{0.6cm}
            {\color{black}\fontsize{24pt}{28pt}\selectfont \CourseName\\
                \CourseCode}

            \vfill
            {\color{black}\fontsize{16pt}{18pt}\selectfont \Faculty \par \Institution}

            \vfill
            {\color{black}\fontsize{14pt}{18pt}\selectfont \Instructor \par \Email}

            \vspace{0.6cm}
            \color{BrandGreen}\rule{0.9\textwidth}{0.5pt}
        \end{frame}
    }

% =========================================================
% ESPACIADO AUTOMÁTICO ANTES DE ECUACIONES DISPLAY
% =========================================================

\newlength{\DisplayMathTopSpace}
\setlength{\DisplayMathTopSpace}{-0.2cm}

\let\OriginalDisplayMath\[
\renewcommand{\[}{%
    \vspace{\DisplayMathTopSpace}%
    \OriginalDisplayMath
}


% =========================================================
% CONFIGURACIÓN DEL CURSO
% =========================================================
\renewcommand{\CourseName}{Robótica del Servicio}
\renewcommand{\CourseCode}{ING 01335}
\renewcommand{\Faculty}{Facultad de Ingeniería}
\renewcommand{\Institution}{Politécnico Colombiano Jaime Isaza Cadavid}
\renewcommand{\Instructor}{Deimer Miranda Montoya, MSc.(c).}
\renewcommand{\Email}{deimer\_miranda91162@elpoli.edu.co}
\renewcommand{\InstitutionLogo}{Template/LogoPoli.png}

\title{Robótica del Servicio}
\subtitle{Clase 2: Geometría en el plano, Linux y Python}
\author{\Instructor}
\institute{\Institution\\\Faculty}
\date{2026-2}

\begin{document}

\TitleSlide
\NormasSlide

% =========================================================
% CLASE 3
% =========================================================
\section{Clase 3}

\SectionSlide
{Clase 3: Rotaciones y transformaciones en el plano}
{Del sistema local del robot al sistema global}

% =========================================================
\begin{frame}{Propósito de la clase}

    Al finalizar esta sesión, el estudiante estará en capacidad de:

    \begin{itemize}
        \item interpretar geométricamente una rotación en el plano;
        \item identificar el sentido positivo y negativo de una rotación;
        \item deducir la matriz de rotación bidimensional;
        \item aplicar matrices de rotación a puntos y vectores.
    \end{itemize}

    \vspace{0.3cm}

    \Idea{
        La rotación permite describir cómo cambia la orientación de un punto,
        vector o sistema de referencia en el plano.
    }

\end{frame}

% =========================================================
\begin{frame}{Propósito de la clase}

    Además, el estudiante estará en capacidad de:

    \begin{itemize}
        \item comprender las principales propiedades de $R(\theta)$;
        \item realizar composiciones e inversiones de rotaciones;
        \item transformar coordenadas entre el sistema local del robot y el sistema global;
        \item validar y representar los resultados utilizando Python, NumPy y Matplotlib.
    \end{itemize}

    \vspace{0.3cm}

    \Idea{
        La matriz de rotación permite expresar matemáticamente cómo cambia la
        representación de un punto o vector cuando cambia su orientación.
    }

\end{frame}

% =========================================================
\section{Representaciones de un punto}

\SectionSlide
{Representaciones de un punto}
{Cartesiano y polar}

% =========================================================
\begin{frame}{Coordenadas cartesianas}

    \begin{columns}[T]

        \begin{column}{0.52\textwidth}

            En el plano, un punto puede representarse mediante sus coordenadas:

            \[
            P=(x,y)
            \]

            o en forma vectorial:

            \[
            \mathbf{p}
            =
            \begin{bmatrix}
                x\\
                y
            \end{bmatrix}
            \]

            donde:

            \begin{itemize}
                \item $x$ es la componente horizontal;
                \item $y$ es la componente vertical.
            \end{itemize}

        \end{column}

        \begin{column}{0.45\textwidth}

            \begin{center}
                \begin{tikzpicture}[scale=0.65]

                    \draw[-{Latex[length=3mm]}, thick]
                    (0,0) -- (5.8,0) node[right] {$X$};

                    \draw[-{Latex[length=3mm]}, thick]
                    (0,0) -- (0,4.8) node[above] {$Y$};

                    \coordinate (P) at (4,3);

                    \draw[dashed] (P) -- (4,0);
                    \draw[dashed] (P) -- (0,3);

                    \filldraw[BrandBlue] (P) circle (2.5pt);
                    \node[above right] at (P) {$P=(x,y)$};

                    \node[below] at (4,0) {$x$};
                    \node[left] at (0,3) {$y$};

                \end{tikzpicture}
            \end{center}

        \end{column}

    \end{columns}

\end{frame}

% =========================================================
\begin{frame}{Coordenadas cartesianas}

    \DefBlock[Idea clave]{
        Las coordenadas cartesianas describen un punto a partir de sus
        proyecciones sobre los ejes $X$ y $Y$.
    }

    \vspace{0.3cm}

    \Idea{
        Una coordenada no tiene sentido completo si no sabemos respecto a qué
        sistema de referencia fue medida.
    }

\end{frame}

% =========================================================
\begin{frame}{Coordenadas polares}

    \begin{columns}[T]

        \begin{column}{0.55\textwidth}

            El mismo punto también puede describirse mediante:

            \[
            P=(\rho,\alpha)
            \]

            donde:

            \begin{itemize}
                \item $\rho$ es la distancia desde el origen;
                \item $\alpha$ es el ángulo medido desde el eje positivo $X$.
            \end{itemize}

            \[
            \rho=\sqrt{x^2+y^2}
            \]

            \[
            \alpha=\operatorname{atan2}(y,x)
            \]

        \end{column}

        \begin{column}{0.42\textwidth}

            \begin{center}
                \begin{tikzpicture}[scale=0.65]

                    \draw[-{Latex[length=3mm]}, thick]
                    (0,0) -- (5.8,0) node[right] {$X$};

                    \draw[-{Latex[length=3mm]}, thick]
                    (0,0) -- (0,4.8) node[above] {$Y$};

                    \coordinate (P) at (4,3);

                    \draw[-{Latex[length=3mm]}, very thick, BrandBlue]
                    (0,0) -- (P);

                    \filldraw[BrandYellow] (P) circle (2.5pt);
                    \node[above right] at (P) {$P$};

                    \draw[->, BrandGreen, thick]
                    (1.0,0) arc[start angle=0,end angle=36.87,radius=1.0];

                    \node at (1.25,0.45) {$\alpha$};
                    \node at (2.25,1.9) {$\rho$};

                \end{tikzpicture}
            \end{center}

        \end{column}

    \end{columns}

\end{frame}

% =========================================================
\begin{frame}{Coordenadas polares}

    \Idea{
        En coordenadas polares, un punto queda definido por
        \textbf{magnitud y dirección}.
    }

    \vspace{0.3cm}

    \DefBlock[Comparación]{
        En coordenadas cartesianas usamos componentes $(x,y)$.
        En coordenadas polares usamos distancia y dirección $(\rho,\alpha)$.
    }

\end{frame}

% =========================================================
\begin{frame}{Un mismo punto, dos descripciones}

    \begin{columns}[T]

        \begin{column}{0.47\textwidth}

            \[
            P=(x,y)
            \]

            describe el punto mediante sus componentes sobre los ejes.

            \vspace{0.3cm}

            \[
            P=(\rho,\alpha)
            \]

            describe el mismo punto mediante su distancia al origen y su dirección.

        \end{column}

        \begin{column}{0.50\textwidth}

            \begin{center}
                \begin{tikzpicture}[scale=0.62]

                    \draw[-{Latex[length=3mm]}, thick]
                    (0,0) -- (5.8,0) node[right] {$X$};

                    \draw[-{Latex[length=3mm]}, thick]
                    (0,0) -- (0,4.8) node[above] {$Y$};

                    \coordinate (P) at (4,3);

                    \draw[-{Latex[length=3mm]}, very thick, BrandBlue]
                    (0,0) -- (P);

                    \draw[dashed] (P) -- (4,0);
                    \draw[dashed] (P) -- (0,3);

                    \filldraw[BrandYellow] (P) circle (2.5pt);

                    \node[above right] at (P) {$P$};
                    \node[below] at (4,0) {$x$};
                    \node[left] at (0,3) {$y$};
                    \node at (2.2,1.9) {$\rho$};

                    \draw[->, BrandGreen, thick]
                    (1.0,0) arc[start angle=0,end angle=36.87,radius=1.0];

                    \node at (1.25,0.45) {$\alpha$};

                \end{tikzpicture}
            \end{center}

        \end{column}

    \end{columns}



    \DefBlock[Equivalencia]{
        Ambas representaciones describen exactamente el mismo punto.
        Lo que cambia es la manera de expresar su posición.
    }

\end{frame}

% =========================================================
\begin{frame}{De polares a cartesianas}

    Si se conocen la magnitud $\rho$ y el ángulo $\alpha$, pueden obtenerse
    las coordenadas cartesianas mediante trigonometría:

    \[
    x=\rho\cos\alpha
    \]

    \[
    y=\rho\sin\alpha
    \]

    por tanto:

    \[
    \boxed{
        \mathbf{p}
        =
        \begin{bmatrix}
            \rho\cos\alpha\\
            \rho\sin\alpha
    \end{bmatrix}}
    \]

\end{frame}

% =========================================================
\begin{frame}{De polares a cartesianas}

    \begin{center}
        \begin{tikzpicture}[scale=0.72]

            \draw[-{Latex[length=3mm]}, thick]
            (0,0) -- (6,0) node[right] {$X$};

            \draw[-{Latex[length=3mm]}, thick]
            (0,0) -- (0,4.8) node[above] {$Y$};

            \coordinate (P) at (4,3);

            \draw[-{Latex[length=3mm]}, very thick, BrandBlue]
            (0,0) -- (P);

            \draw[dashed] (P) -- (4,0);
            \draw[dashed] (P) -- (0,3);

            \filldraw[BrandYellow] (P) circle (2.5pt);

            \node[below] at (2,0) {$x=\rho\cos\alpha$};
            \node[left] at (0,1.5) {$y=\rho\sin\alpha$};

            \node at (2.25,1.9) {$\rho$};

            \draw[->, BrandGreen, thick]
            (1.0,0) arc[start angle=0,end angle=36.87,radius=1.0];

            \node at (1.25,0.45) {$\alpha$};

        \end{tikzpicture}
    \end{center}

    \Result{
        La representación polar permite expresar un punto a partir de una magnitud
        y una dirección.
    }

\end{frame}

% =========================================================
\section{Rotaciones en el plano}

\SectionSlide
{Rotaciones en el plano}
{¿Qué significa girar un vector?}

% =========================================================
\begin{frame}{¿Qué significa rotar un vector?}

    Si el vector se escribe como:

    \[
    \mathbf{p}
    =
    \begin{bmatrix}
        \rho\cos\alpha\\
        \rho\sin\alpha
    \end{bmatrix}
    \]

    y luego se rota un ángulo $\theta$, entonces:

    \[
    \alpha'=\alpha+\theta
    \]

    mientras que la magnitud permanece igual:

    \[
    \rho'=\rho
    \]

\end{frame}

% =========================================================
\begin{frame}{¿Qué significa rotar un vector?}

    Por tanto:

    \[
    \mathbf{p}'
    =
    \begin{bmatrix}
        \rho\cos(\alpha+\theta)\\
        \rho\sin(\alpha+\theta)
    \end{bmatrix}
    \]



    \begin{center}
        \begin{tikzpicture}[scale=0.6]

            \draw[-{Latex[length=3mm]}, thick]
            (0,0) -- (6,0) node[right] {$X$};

            \draw[-{Latex[length=3mm]}, thick]
            (0,0) -- (0,5) node[above] {$Y$};

            \draw[-{Latex[length=3mm]}, very thick, BrandBlue]
            (0,0) -- (4,1.6);

            \draw[-{Latex[length=3mm]}, very thick, BrandGreen]
            (0,0) -- (2.4,3.6);

            \node[right] at (4,1.6) {$\mathbf{p}$};
            \node[above] at (2.4,3.6) {$\mathbf{p}'$};

            \draw[->, BrandYellow, thick]
            (1.2,0.48) arc[start angle=21.8,end angle=56.3,radius=1.3];

            \node at (1.75,1.2) {$\theta$};

        \end{tikzpicture}
    \end{center}

    \Idea{
        Esta idea será el punto de partida para deducir la matriz de rotación
        $R(\theta)$.
    }

\end{frame}

% =========================================================
\begin{frame}{Convención de signos}

    Utilizaremos la convención matemática habitual:

    \begin{itemize}
        \item rotación \textbf{antihoraria}: $\theta>0$;
        \item rotación \textbf{horaria}: $\theta<0$.
    \end{itemize}

    \vspace{0.3cm}

    \begin{center}
        \begin{tikzpicture}[scale=0.85]

            \draw[-{Latex[length=3mm]}, thick]
            (-0.4,0) -- (5,0) node[right] {$X$};

            \draw[-{Latex[length=3mm]}, thick]
            (0,-0.4) -- (0,3.5) node[above] {$Y$};

            \draw[-{Latex[length=3mm]}, very thick, BrandBlue]
            (0,0) -- (3.2,0.7)
            node[right] {$\mathbf{p}$};

            \draw[-{Latex[length=3mm]}, very thick, BrandGreen]
            (0,0) -- (2.2,2.5)
            node[above] {$\mathbf{p}'$};

            \draw[->, thick, BrandYellow]
            (1.1,0.24)
            arc[start angle=12,end angle=49,radius=1.12];

            \node at (1.45,0.85) {$+\theta$};

        \end{tikzpicture}
    \end{center}

\end{frame}

% =========================================================
\section{Deducción de la matriz de rotación}

\SectionSlide
{Matriz de rotación}
{Construcción a partir de la trigonometría}

% =========================================================
\begin{frame}{Partimos de la geometría}

    Antes de la rotación:

    \[
    x=\rho\cos\alpha
    \]

    \[
    y=\rho\sin\alpha
    \]

    Después de una rotación $\theta$:

    \[
    x'=\rho\cos(\alpha+\theta)
    \]

    \[
    y'=\rho\sin(\alpha+\theta)
    \]

\end{frame}

% =========================================================
\begin{frame}{Identidades trigonométricas}

    Ahora utilizamos las identidades trigonométricas:

    \[
    \cos(\alpha+\theta)
    =
    \cos\alpha\cos\theta
    -
    \sin\alpha\sin\theta
    \]

    \[
    \sin(\alpha+\theta)
    =
    \sin\alpha\cos\theta
    +
    \cos\alpha\sin\theta
    \]

    \vspace{0.3cm}

    \Idea{
        Estas identidades permiten pasar de la interpretación geométrica
        a una expresión matricial.
    }

\end{frame}

% =========================================================
\begin{frame}{Desarrollo de la coordenada \texorpdfstring{$x'$}{x'}}

    Partimos de:

    \[
    x'=\rho\cos(\alpha+\theta)
    \]

    Sustituyendo la identidad trigonométrica:

    \[
    x'
    =
    \rho
    \left(
    \cos\alpha\cos\theta
    -
    \sin\alpha\sin\theta
    \right)
    \]

    Distribuyendo $\rho$:

    \[
    x'
    =
    \rho\cos\alpha\cos\theta
    -
    \rho\sin\alpha\sin\theta
    \]

\end{frame}

% =========================================================
\begin{frame}{Desarrollo de la coordenada \texorpdfstring{$x'$}{x'}}

    Sabemos que:

    \[
    \rho\cos\alpha=x,
    \qquad
    \rho\sin\alpha=y
    \]

    por tanto:

    \[
    x'
    =
    x\cos\theta
    -
    y\sin\theta
    \]

    \[
    \boxed{
        x'=x\cos\theta-y\sin\theta
    }
    \]

\end{frame}

% =========================================================
\begin{frame}{Desarrollo de la coordenada \texorpdfstring{$y'$}{y'}}

    Partimos de:

    \[
    y'=\rho\sin(\alpha+\theta)
    \]

    Sustituyendo:

    \[
    y'
    =
    \rho
    \left(
    \sin\alpha\cos\theta
    +
    \cos\alpha\sin\theta
    \right)
    \]

    Distribuyendo:

    \[
    y'
    =
    \rho\sin\alpha\cos\theta
    +
    \rho\cos\alpha\sin\theta
    \]

\end{frame}

% =========================================================
\begin{frame}{Desarrollo de la coordenada \texorpdfstring{$y'$}{y'}}

    Sabemos que:

    \[
    \rho\cos\alpha=x,
    \qquad
    \rho\sin\alpha=y
    \]

    por tanto:

    \[
    y'
    =
    y\cos\theta
    +
    x\sin\theta
    \]

    \[
    \boxed{
        y'=x\sin\theta+y\cos\theta
    }
    \]

\end{frame}

% =========================================================
\begin{frame}{Escribamos las dos ecuaciones}

    Los resultados anteriores son:

    \[
    x'
    =
    x\cos\theta-y\sin\theta
    \]

    \[
    y'
    =
    x\sin\theta+y\cos\theta
    \]

    \vspace{0.3cm}

    Estas dos ecuaciones describen cómo cambian las coordenadas de un vector
    cuando se rota un ángulo $\theta$.

\end{frame}

% =========================================================
\begin{frame}{Forma matricial}

    Podemos organizarlas matricialmente:

    \[
    \begin{bmatrix}
        x'\\
        y'
    \end{bmatrix}
    =
    \begin{bmatrix}
        \cos\theta & -\sin\theta\\
        \sin\theta & \cos\theta
    \end{bmatrix}
    \begin{bmatrix}
        x\\
        y
    \end{bmatrix}
    \]

    Definimos:

    \[
    \boxed{
        R(\theta)
        =
        \begin{bmatrix}
            \cos\theta & -\sin\theta\\
            \sin\theta & \cos\theta
    \end{bmatrix}}
    \]

\end{frame}

% =========================================================
\begin{frame}{Matriz de rotación}

    La operación de rotación puede escribirse de manera compacta:

    \[
    \boxed{
        \mathbf{p}'=R(\theta)\mathbf{p}
    }
    \]

    donde:

    \[
    R(\theta)
    =
    \begin{bmatrix}
        \cos\theta & -\sin\theta\\
        \sin\theta & \cos\theta
    \end{bmatrix}
    \]

    y:

    \[
    \mathbf{p}
    =
    \begin{bmatrix}
        x\\
        y
    \end{bmatrix}
    \]

\end{frame}

% =========================================================
\begin{frame}{Matriz de rotación}

    \DefBlock[Matriz de rotación]{
        Es una transformación lineal que modifica la orientación de un vector
        sin modificar su magnitud.
    }

    \vspace{0.3cm}

    \Idea{
        Esta matriz será la base para relacionar el sistema local del robot
        con el sistema global.
    }

\end{frame}

% =========================================================
\section{Aplicación de la matriz}

\SectionSlide
{Aplicación de la matriz de rotación}
{Del modelo matemático al cálculo}

% =========================================================
\begin{frame}{Ejemplo: rotación de \texorpdfstring{$90^\circ$}{90 grados}}

    Considere:

    \[
    \mathbf{p}
    =
    \begin{bmatrix}
        3\\
        1
    \end{bmatrix}
    \]

    y una rotación:

    \[
    \theta=90^\circ
    =
    \frac{\pi}{2}
    \]

    La matriz será:

    \[
    R\left(\frac{\pi}{2}\right)
    =
    \begin{bmatrix}
        \cos\frac{\pi}{2} & -\sin\frac{\pi}{2}\\
        \sin\frac{\pi}{2} & \cos\frac{\pi}{2}
    \end{bmatrix}
    \]

\end{frame}

% =========================================================
\begin{frame}{Ejemplo: rotación de \texorpdfstring{$90^\circ$}{90 grados}}

    Como:

    \[
    \cos\frac{\pi}{2}=0,
    \qquad
    \sin\frac{\pi}{2}=1
    \]

    entonces:

    \[
    R\left(\frac{\pi}{2}\right)
    =
    \begin{bmatrix}
        0 & -1\\
        1 & 0
    \end{bmatrix}
    \]

\end{frame}

% =========================================================
\begin{frame}{Ejemplo: operación matricial}

    Aplicamos:

    \[
    \mathbf{p}'
    =
    R\left(\frac{\pi}{2}\right)\mathbf{p}
    \]

    \[
    \mathbf{p}'
    =
    \begin{bmatrix}
        0 & -1\\
        1 & 0
    \end{bmatrix}
    \begin{bmatrix}
        3\\
        1
    \end{bmatrix}
    \]

    Realizando el producto:

    \[
    x'
    =
    (0)(3)+(-1)(1)
    =
    -1
    \]

    \[
    y'
    =
    (1)(3)+(0)(1)
    =
    3
    \]

\end{frame}

% =========================================================
\begin{frame}{Ejemplo: operación matricial}

    Por tanto:

    \[
    \boxed{
        \mathbf{p}'
        =
        \begin{bmatrix}
            -1\\
            3
    \end{bmatrix}}
    \]

    \vspace{0.4cm}

    \Result{
        El vector $(3,1)$ rotado $90^\circ$ en sentido antihorario
        se convierte en $(-1,3)$.
    }

\end{frame}

% =========================================================
\begin{frame}{Comprobación de la magnitud}

    Antes de rotar:

    \[
    \|\mathbf{p}\|
    =
    \sqrt{3^2+1^2}
    =
    \sqrt{10}
    \]

    Después de rotar:

    \[
    \|\mathbf{p}'\|
    =
    \sqrt{(-1)^2+3^2}
    =
    \sqrt{10}
    \]

    Por tanto:

    \[
    \boxed{
        \|\mathbf{p}'\|=\|\mathbf{p}\|
    }
    \]

\end{frame}

% =========================================================
\begin{frame}{Comprobación de la magnitud}

    \Idea{
        Una rotación modifica las coordenadas del vector, pero conserva
        su longitud.
    }

    \vspace{0.3cm}

    \Result{
        Esta propiedad será muy importante cuando representemos velocidades,
        mediciones y sistemas de referencia.
    }

\end{frame}

% =========================================================
\begin{frame}{Otro ejemplo: rotación de \texorpdfstring{$-45^\circ$}{-45 grados}}

    Considere:

    \[
    \mathbf{p}
    =
    \begin{bmatrix}
        2\\
        4
    \end{bmatrix},
    \qquad
    \theta=-45^\circ
    \]

    Primero:

    \[
    \theta
    =
    -\frac{\pi}{4}
    \]

    Entonces:

    \[
    R\left(-\frac{\pi}{4}\right)
    =
    \begin{bmatrix}
        \cos(-\pi/4) & -\sin(-\pi/4)\\
        \sin(-\pi/4) & \cos(-\pi/4)
    \end{bmatrix}
    \]

\end{frame}

% =========================================================
\begin{frame}{Otro ejemplo: rotación de \texorpdfstring{$-45^\circ$}{-45 grados}}

    Utilizando:

    \[
    \cos(-\pi/4)=\frac{\sqrt{2}}{2}
    \]

    \[
    \sin(-\pi/4)=-\frac{\sqrt{2}}{2}
    \]

    La matriz resulta:

    \[
    R\left(-\frac{\pi}{4}\right)
    =
    \begin{bmatrix}
        \frac{\sqrt{2}}{2} & \frac{\sqrt{2}}{2}\\
        -\frac{\sqrt{2}}{2} & \frac{\sqrt{2}}{2}
    \end{bmatrix}
    \]

\end{frame}

% =========================================================
\begin{frame}{Rotación de \texorpdfstring{$-45^\circ$}{-45 grados}: cálculo}

    Aplicando la rotación:

    \[
    \mathbf{p}'
    =
    R\left(-\frac{\pi}{4}\right)
    \begin{bmatrix}
        2\\
        4
    \end{bmatrix}
    \]

    \[
    x'
    =
    \frac{\sqrt{2}}{2}(2)
    +
    \frac{\sqrt{2}}{2}(4)
    =
    3\sqrt{2}
    \]

    \[
    y'
    =
    -\frac{\sqrt{2}}{2}(2)
    +
    \frac{\sqrt{2}}{2}(4)
    =
    \sqrt{2}
    \]

\end{frame}

% =========================================================
\begin{frame}{Rotación de \texorpdfstring{$-45^\circ$}{-45 grados}: resultado}

    Por tanto:

    \[
    \boxed{
        \mathbf{p}'
        =
        \begin{bmatrix}
            3\sqrt{2}\\
            \sqrt{2}
    \end{bmatrix}}
    \]

    \vspace{0.4cm}

    \Result{
        Este ejemplo muestra una rotación en sentido horario,
        por eso el ángulo utilizado es negativo.
    }

\end{frame}

% =========================================================
\section{Propiedades de la matriz de rotación}

\SectionSlide
{Propiedades de $R(\theta)$}
{¿Qué hace especial a una matriz de rotación?}

% =========================================================
\begin{frame}{Rotación nula}

    Para:

    \[
    \theta=0
    \]

    tenemos:

    \[
    R(0)
    =
    \begin{bmatrix}
        \cos0 & -\sin0\\
        \sin0 & \cos0
    \end{bmatrix}
    \]

    Como:

    \[
    \cos0=1,
    \qquad
    \sin0=0
    \]

\end{frame}

% =========================================================
\begin{frame}{Rotación nula}

    Resulta:

    \[
    \boxed{
        R(0)
        =
        \begin{bmatrix}
            1&0\\
            0&1
        \end{bmatrix}
        =
        I
    }
    \]

    Por tanto:

    \[
    R(0)\mathbf{p}=\mathbf{p}
    \]

    \vspace{0.3cm}

    \Idea{
        La rotación de cero grados no modifica el vector.
    }

\end{frame}

% =========================================================
\begin{frame}{Rotación inversa}

    Si primero realizamos una rotación $\theta$, para regresar a la
    orientación original debemos aplicar:

    \[
    -\theta
    \]

    Por tanto:

    \[
    R^{-1}(\theta)=R(-\theta)
    \]

    Además:

    \[
    R(\theta)R(-\theta)=I
    \]

    \vspace{0.3cm}

    \DefBlock[Interpretación]{
        La matriz inversa permite deshacer una rotación.
    }

\end{frame}

% =========================================================
\begin{frame}{Transpuesta e inversa}

    La transpuesta de:

    \[
    R(\theta)
    =
    \begin{bmatrix}
        \cos\theta & -\sin\theta\\
        \sin\theta & \cos\theta
    \end{bmatrix}
    \]

    es:

    \[
    R^T(\theta)
    =
    \begin{bmatrix}
        \cos\theta & \sin\theta\\
        -\sin\theta & \cos\theta
    \end{bmatrix}
    \]

\end{frame}

% =========================================================
\begin{frame}{Transpuesta e inversa}

    Observe que:

    \[
    R^T(\theta)=R(-\theta)
    \]

    y por tanto:

    \[
    \boxed{
        R^{-1}(\theta)=R^T(\theta)
    }
    \]

    \vspace{0.4cm}

    \Result{
        Las matrices de rotación son matrices ortogonales.
    }

\end{frame}

% =========================================================
\begin{frame}{Ortogonalidad}

    Una matriz es ortogonal cuando:

    \[
    R^TR=I
    \]

    Para una matriz de rotación:

    \[
    R^T(\theta)R(\theta)
    =
    I
    \]

    Esta propiedad explica matemáticamente por qué se conserva la magnitud.

    \[
    \mathbf{p}'=R\mathbf{p}
    \]

\end{frame}

% =========================================================
\begin{frame}{Ortogonalidad y magnitud}

    Si:

    \[
    \mathbf{p}'=R\mathbf{p}
    \]

    entonces:

    \[
    \|\mathbf{p}'\|^2
    =
    \mathbf{p}'^T\mathbf{p}'
    \]

    \[
    =
    (R\mathbf{p})^T(R\mathbf{p})
    \]

    \[
    =
    \mathbf{p}^TR^TR\mathbf{p}
    \]

    \[
    =
    \mathbf{p}^T\mathbf{p}
    \]

    por tanto:

    \[
    \boxed{\|\mathbf{p}'\|=\|\mathbf{p}\|}
    \]

\end{frame}

% =========================================================
\begin{frame}{Determinante de una rotación}

    Para:

    \[
    R(\theta)
    =
    \begin{bmatrix}
        \cos\theta & -\sin\theta\\
        \sin\theta & \cos\theta
    \end{bmatrix}
    \]

    el determinante es:

    \[
    \det(R)
    =
    \cos^2\theta+\sin^2\theta
    \]

    utilizando:

    \[
    \cos^2\theta+\sin^2\theta=1
    \]

    obtenemos:

    \[
    \boxed{
        \det(R)=1
    }
    \]

\end{frame}

% =========================================================
\begin{frame}{Determinante de una rotación}

    \Idea{
        Una rotación pura conserva longitudes, ángulos y orientación geométrica
        del sistema.
    }

    \vspace{0.3cm}

    \Result{
        El determinante igual a uno indica que la transformación no escala ni
        refleja el espacio.
    }

\end{frame}

% =========================================================
\section{Composición de rotaciones}

\SectionSlide
{Composición de rotaciones}
{¿Qué ocurre cuando giramos varias veces?}

% =========================================================
\begin{frame}{Dos rotaciones consecutivas}

    Supongamos que primero aplicamos una rotación $\theta_1$:

    \[
    \mathbf{p}_1
    =
    R(\theta_1)\mathbf{p}
    \]

    y posteriormente una rotación $\theta_2$:

    \[
    \mathbf{p}_2
    =
    R(\theta_2)\mathbf{p}_1
    \]

    Sustituyendo:

    \[
    \mathbf{p}_2
    =
    R(\theta_2)
    R(\theta_1)
    \mathbf{p}
    \]

\end{frame}

% =========================================================
\begin{frame}{Dos rotaciones consecutivas}

    La rotación equivalente es:

    \[
    \boxed{
        R(\theta_2)R(\theta_1)
        =
        R(\theta_1+\theta_2)
    }
    \]

    \vspace{0.3cm}

    \Idea{
        Varias rotaciones alrededor del mismo origen pueden representarse como
        una única rotación equivalente.
    }

\end{frame}

% =========================================================
\begin{frame}{Ejemplo de composición}

    Supongamos:

    \[
    \theta_1=30^\circ
    \]

    \[
    \theta_2=45^\circ
    \]

    La rotación equivalente será:

    \[
    \theta_T
    =
    30^\circ+45^\circ
    \]

    \[
    \boxed{
        \theta_T=75^\circ
    }
    \]

\end{frame}

% =========================================================
\begin{frame}{Ejemplo de composición}

    Por tanto:

    \[
    R(45^\circ)R(30^\circ)
    =
    R(75^\circ)
    \]

    \vspace{0.3cm}

    \Result{
        Varias rotaciones alrededor del mismo origen pueden combinarse
        en una única rotación equivalente.
    }

\end{frame}

% =========================================================
\section{Sistemas de referencia}

\SectionSlide
{Sistemas de referencia}
{¿Cómo ve el robot un punto y cómo lo ve el mundo?}

% =========================================================
\begin{frame}{Recordemos los dos sistemas}

    Consideraremos dos sistemas de referencia:

    \begin{itemize}
        \item $\{G\}$: sistema global;
        \item $\{R\}$: sistema local unido al robot.
    \end{itemize}

    La pose del robot respecto al sistema global es:

    \[
    {}^{G}\mathbf{x}_R
    =
    \begin{bmatrix}
        x_R\\
        y_R\\
        \theta
    \end{bmatrix}
    \]

\end{frame}

% =========================================================
\begin{frame}{Recordemos los dos sistemas}

    \Idea{
        El robot puede medir objetos respecto a sus propios ejes,
        pero para navegar necesitamos frecuentemente expresar esas posiciones
        respecto al mundo.
    }

    \vspace{0.3cm}

    \DefBlock[Sistema global y sistema local]{
        El sistema global describe el mundo. El sistema local se mueve
        con el robot y depende de su posición y orientación.
    }

\end{frame}

% =========================================================
\begin{frame}{Un punto expresado en el robot}

    Supongamos que un sensor detecta un objeto en:

    \[
    {}^{R}\mathbf{p}
    =
    \begin{bmatrix}
        x_R^{p}\\
        y_R^{p}
    \end{bmatrix}
    \]

    La notación:

    \[
    {}^{R}\mathbf{p}
    \]

    significa:

    \begin{center}
        ``las coordenadas del punto $\mathbf{p}$ están expresadas
        respecto al sistema del robot $\{R\}$''.
    \end{center}

\end{frame}

% =========================================================
\begin{frame}{Un punto expresado en el robot}

    \DefBlock[Importante]{
        El mismo punto físico puede tener coordenadas diferentes dependiendo
        del sistema de referencia desde el cual se describa.
    }

    \vspace{0.3cm}

    \Idea{
        Por eso siempre debemos indicar respecto a qué sistema de referencia
        están expresadas las coordenadas.
    }

\end{frame}

% =========================================================
\begin{frame}{Primero: considerar únicamente orientación}

    Supongamos inicialmente que los orígenes de ambos sistemas coinciden.

    Si el sistema del robot está rotado un ángulo $\theta$ respecto al global:

    \[
    {}^{G}\mathbf{p}
    =
    R(\theta)
    {}^{R}\mathbf{p}
    \]

    es decir:

    \[
    \boxed{
        {}^{G}\mathbf{p}
        =
        \begin{bmatrix}
            \cos\theta & -\sin\theta\\
            \sin\theta & \cos\theta
        \end{bmatrix}
        {}^{R}\mathbf{p}
    }
    \]

\end{frame}

% =========================================================
\begin{frame}{Primero: considerar únicamente orientación}

    \Idea{
        La matriz $R(\theta)$ permite expresar en coordenadas globales
        un vector que originalmente estaba expresado en coordenadas del robot.
    }

    \vspace{0.3cm}

    \Result{
        Si solo cambia la orientación, basta con aplicar una rotación.
    }

\end{frame}

% =========================================================
\begin{frame}{Pero el robot también está trasladado}

    En un escenario real, el robot normalmente no se encuentra en el origen.

    Su posición global es:

    \[
    {}^{G}\mathbf{p}_R
    =
    \begin{bmatrix}
        x_R\\
        y_R
    \end{bmatrix}
    \]

    Por tanto, para transformar un punto local al sistema global debemos:

    \begin{enumerate}
        \item rotar el punto según la orientación del robot;
        \item trasladarlo hasta la posición global del robot.
    \end{enumerate}

\end{frame}

% =========================================================
\begin{frame}{Pero el robot también está trasladado}

    Esto produce:

    \[
    \boxed{
        {}^{G}\mathbf{p}
        =
        {}^{G}\mathbf{p}_R
        +
        R(\theta)
        {}^{R}\mathbf{p}
    }
    \]

    \vspace{0.3cm}

    \Idea{
        Esta ecuación combina dos operaciones: rotación y traslación.
    }

\end{frame}

% =========================================================
\begin{frame}{Interpretación de la transformación}

    La expresión:

    \[
    {}^{G}\mathbf{p}
    =
    {}^{G}\mathbf{p}_R
    +
    R(\theta)
    {}^{R}\mathbf{p}
    \]

    contiene dos operaciones diferentes:

    \vspace{0.3cm}

    \[
    \underbrace{
        R(\theta)\,{}^{R}\mathbf{p}
    }_{\text{rotación}}
    \]

    y:

    \[
    \underbrace{
        {}^{G}\mathbf{p}_R
    }_{\text{traslación}}
    \]

\end{frame}

% =========================================================
\begin{frame}{Interpretación de la transformación}

    \Result{
        Rotación y traslación permiten relacionar matemáticamente
        el sistema del robot con el sistema global.
    }

    \vspace{0.3cm}

    \Idea{
        Esta operación será fundamental para interpretar mediciones de sensores
        y ubicarlas dentro del mundo.
    }

\end{frame}

% =========================================================
\begin{frame}{Ejemplo: punto local a coordenadas globales}

    Suponga que el robot está en:

    \[
    {}^{G}\mathbf{p}_R
    =
    \begin{bmatrix}
        2\\
        3
    \end{bmatrix}
    \]

    con orientación:

    \[
    \theta=90^\circ
    \]

    y detecta un objeto a:

    \[
    {}^{R}\mathbf{p}
    =
    \begin{bmatrix}
        2\\
        0
    \end{bmatrix}
    \]

    respecto a su propio sistema.

    Queremos calcular:

    \[
    {}^{G}\mathbf{p}
    \]

\end{frame}

% =========================================================
\begin{frame}{Ejemplo: aplicar la rotación}

    Para:

    \[
    \theta=90^\circ
    \]

    tenemos:

    \[
    R(90^\circ)
    =
    \begin{bmatrix}
        0&-1\\
        1&0
    \end{bmatrix}
    \]

    Entonces:

    \[
    R(90^\circ)
    {}^{R}\mathbf{p}
    =
    \begin{bmatrix}
        0&-1\\
        1&0
    \end{bmatrix}
    \begin{bmatrix}
        2\\
        0
    \end{bmatrix}
    \]

    \[
    =
    \begin{bmatrix}
        0\\
        2
    \end{bmatrix}
    \]

\end{frame}

% =========================================================
\begin{frame}{Ejemplo: aplicar la rotación}

    \Result{
        El objeto se encuentra $2$ unidades en la dirección global $Y$
        respecto al robot.
    }

    \vspace{0.3cm}

    \Idea{
        La rotación cambió la forma en que se expresa el punto local,
        pero todavía falta sumar la posición global del robot.
    }

\end{frame}

% =========================================================
\begin{frame}{Ejemplo: aplicar la traslación}

    Ahora sumamos la posición global del robot:

    \[
    {}^{G}\mathbf{p}
    =
    \begin{bmatrix}
        2\\
        3
    \end{bmatrix}
    +
    \begin{bmatrix}
        0\\
        2
    \end{bmatrix}
    \]

    Por tanto:

    \[
    \boxed{
        {}^{G}\mathbf{p}
        =
        \begin{bmatrix}
            2\\
            5
    \end{bmatrix}}
    \]

\end{frame}

% =========================================================
\begin{frame}{Ejemplo: aplicar la traslación}

    \Result{
        Aunque el objeto se encuentra en $(2,0)$ respecto al robot,
        su posición respecto al sistema global es $(2,5)$.
    }

    \vspace{0.3cm}

    \Idea{
        El punto local primero se rota y luego se traslada hasta la ubicación
        global del robot.
    }

\end{frame}

% =========================================================
\begin{frame}{Representación geométrica}

    \begin{center}
        \begin{tikzpicture}[scale=0.72]

            % Sistema global
            \draw[-{Latex[length=3mm]}, thick]
            (0,0) -- (7,0) node[right] {$X_G$};

            \draw[-{Latex[length=3mm]}, thick]
            (0,0) -- (0,6) node[above] {$Y_G$};

            % Robot
            \coordinate (R) at (2,3);

            \begin{scope}[shift={(R)},rotate=90]

                \draw[
                fill=BrandBlue!15,
                draw=BrandBlue,
                thick
                ]
                (-0.4,-0.3) rectangle (0.4,0.3);

                \draw[-{Latex[length=2.5mm]}, very thick, BrandGreen]
                (0,0) -- (1.4,0) node[right] {$x_R$};

                \draw[-{Latex[length=2.5mm]}, very thick, BrandGreen]
                (0,0) -- (0,1.2) node[above] {$y_R$};

            \end{scope}

            \node[right] at (2.15,2.7) {Robot};

            % Punto global
            \filldraw[BrandYellow]
            (2,5) circle (3pt);

            \node[right] at (2.1,5)
            {$P$};

            % Línea local
            \draw[dashed, thick]
            (R) -- (2,5);

            \node[right] at (2.15,4)
            {$2$};

            % Proyecciones
            \draw[dashed]
            (2,0) -- (2,3);

            \draw[dashed]
            (0,3) -- (2,3);

            \node[below] at (2,0) {$2$};
            \node[left] at (0,3) {$3$};

        \end{tikzpicture}
    \end{center}

\end{frame}

% =========================================================
\section{Transformación inversa}

\SectionSlide
{Transformación inversa}
{Del sistema global al sistema del robot}

% =========================================================
\begin{frame}{¿Y si conocemos el punto global?}

    Partimos de:

    \[
    {}^{G}\mathbf{p}
    =
    {}^{G}\mathbf{p}_R
    +
    R(\theta)
    {}^{R}\mathbf{p}
    \]

    Restamos la posición del robot:

    \[
    {}^{G}\mathbf{p}
    -
    {}^{G}\mathbf{p}_R
    =
    R(\theta)
    {}^{R}\mathbf{p}
    \]

    Multiplicamos por la inversa:

    \[
    {}^{R}\mathbf{p}
    =
    R^{-1}(\theta)
    \left(
    {}^{G}\mathbf{p}
    -
    {}^{G}\mathbf{p}_R
    \right)
    \]

\end{frame}

% =========================================================
\begin{frame}{¿Y si conocemos el punto global?}

    Como:

    \[
    R^{-1}(\theta)=R^T(\theta)
    \]

    obtenemos:

    \[
    \boxed{
        {}^{R}\mathbf{p}
        =
        R^T(\theta)
        \left(
        {}^{G}\mathbf{p}
        -
        {}^{G}\mathbf{p}_R
        \right)}
    \]

    \vspace{0.3cm}

    \Result{
        Esta expresión permite representar un punto global desde el sistema
        local del robot.
    }

\end{frame}

% =========================================================
\begin{frame}{Interpretación de la transformación inversa}

    Para pasar de coordenadas globales a coordenadas del robot:

    \begin{enumerate}
        \item restamos la posición global del robot;
        \item aplicamos la rotación inversa.
    \end{enumerate}

    Matemáticamente:

    \[
    {}^{R}\mathbf{p}
    =
    R(-\theta)
    \left(
    {}^{G}\mathbf{p}
    -
    {}^{G}\mathbf{p}_R
    \right)
    \]

\end{frame}

% =========================================================
\begin{frame}{Interpretación de la transformación inversa}

    \Idea{
        La transformación directa responde:

        \textit{¿Dónde está este punto del robot visto desde el mundo?}

        La transformación inversa responde:

        \textit{¿Dónde está este punto del mundo visto desde el robot?}
    }

\end{frame}

% =========================================================
\section{Validación con Python}

\SectionSlide
{Validación computacional}
{NumPy como herramienta para comprobar la matemática}

% =========================================================
\begin{frame}[fragile]{Construcción de la matriz con NumPy}

    \begin{pythoncodedark}
        import numpy as np

        theta_deg = 60
        theta = np.deg2rad(theta_deg)

        R = np.array([
            [np.cos(theta), -np.sin(theta)],
            [np.sin(theta),  np.cos(theta)]
            ])

        print(R)
    \end{pythoncodedark}

\end{frame}

% =========================================================
\begin{frame}{Construcción de la matriz con NumPy}

    \DefBlock[Importante]{
        Las funciones trigonométricas de NumPy trabajan con ángulos
        expresados en radianes.
    }

    \vspace{0.3cm}

    \Idea{
        Por esta razón, primero convertimos los grados a radianes usando
        \texttt{np.deg2rad}.
    }

\end{frame}

% =========================================================
\begin{frame}[fragile]{Rotar un vector con NumPy}

    \begin{pythoncodedark}
        import numpy as np

        theta = np.deg2rad(60)

        R = np.array([
            [np.cos(theta), -np.sin(theta)],
            [np.sin(theta),  np.cos(theta)]
            ])

        p = np.array([4.0, 1.0])

        p_rotado = R @ p

        print("Vector original:", p)
        print("Vector rotado:", p_rotado)
    \end{pythoncodedark}

\end{frame}

% =========================================================
\begin{frame}{Rotar un vector con NumPy}

    \Idea{
        El operador \texttt{@} realiza el producto matricial.
    }

    \vspace{0.3cm}

    \Result{
        El resultado computacional debe coincidir con el producto
        $R(\theta)\mathbf{p}$ desarrollado manualmente.
    }

\end{frame}

% =========================================================
\begin{frame}[fragile]{Comprobar la conservación de la magnitud}

    \begin{pythoncodedark}
        import numpy as np

        theta = np.deg2rad(60)

        R = np.array([
            [np.cos(theta), -np.sin(theta)],
            [np.sin(theta),  np.cos(theta)]
            ])

        p = np.array([4.0, 1.0])
        p_rotado = R @ p

        print("Norma original:", np.linalg.norm(p))
        print("Norma rotada:", np.linalg.norm(p_rotado))
    \end{pythoncodedark}

\end{frame}

% =========================================================
\begin{frame}{Comprobar la conservación de la magnitud}

    \Result{
        La norma del vector antes y después de la rotación debe ser la misma,
        salvo pequeñas diferencias numéricas por redondeo.
    }

    \vspace{0.3cm}

    \Idea{
        Esta prueba computacional confirma una propiedad matemática:
        una rotación cambia la dirección, pero no la longitud del vector.
    }

\end{frame}

% =========================================================
\begin{frame}[fragile]{Verificar propiedades de \texorpdfstring{$R$}{R}}

    \begin{pythoncodedark}
        import numpy as np

        theta = np.deg2rad(35)

        R = np.array([
            [np.cos(theta), -np.sin(theta)],
            [np.sin(theta),  np.cos(theta)]
            ])

        print("R:")
        print(R)

        print("R.T @ R:")
        print(R.T @ R)

        print("Determinante:")
        print(np.linalg.det(R))
    \end{pythoncodedark}

\end{frame}

% =========================================================
\begin{frame}{Verificar propiedades de \texorpdfstring{$R$}{R}}

    El resultado esperado es aproximadamente:

    \[
    R^TR=I
    \]

    y:

    \[
    \det(R)=1
    \]

    \vspace{0.3cm}

    \Result{
        NumPy permite verificar computacionalmente las propiedades algebraicas
        de la matriz de rotación.
    }

\end{frame}

% =========================================================
\begin{frame}[fragile]{Transformación local a global}

    Suponga:

    \[
    {}^{G}\mathbf{p}_R=(3,2),
    \qquad
    \theta=40^\circ,
    \qquad
    {}^{R}\mathbf{p}=(2,1)
    \]

    \begin{pythoncodedark}
        import numpy as np

        p_robot = np.array([3.0, 2.0])
        p_local = np.array([2.0, 1.0])

        theta = np.deg2rad(40)
    \end{pythoncodedark}

\end{frame}

% =========================================================
\begin{frame}[fragile]{Transformación local a global}

    \begin{pythoncodedark}
        R = np.array([
            [np.cos(theta), -np.sin(theta)],
            [np.sin(theta),  np.cos(theta)]
            ])

        p_global = p_robot + R @ p_local

        print("Punto global:", p_global)
    \end{pythoncodedark}



    \Idea{
        Este código implementa la relación:
        \[
        {}^{G}\mathbf{p}
        =
        {}^{G}\mathbf{p}_R+
        R(\theta){}^{R}\mathbf{p}
        \]
    }

\end{frame}

% =========================================================
\section{Visualización}

\SectionSlide
{Visualización geométrica}
{Verificar que el resultado tenga sentido}

% =========================================================
\begin{frame}[fragile]{Graficar un vector y su rotación}

    \begin{pythoncodedark}
        import numpy as np
        import matplotlib.pyplot as plt

        p = np.array([4.0, 1.0])
        theta = np.deg2rad(60)

        R = np.array([
            [np.cos(theta), -np.sin(theta)],
            [np.sin(theta),  np.cos(theta)]
            ])

        pr = R @ p
    \end{pythoncodedark}

\end{frame}

% =========================================================
\begin{frame}[fragile]{Graficar un vector y su rotación}

    \begin{pythoncodedark}
        plt.quiver(0, 0, p[0], p[1],
        angles="xy", scale_units="xy",
        scale=1, label="Original")

        plt.quiver(0, 0, pr[0], pr[1],
        angles="xy", scale_units="xy",
        scale=1, label="Rotado")
    \end{pythoncodedark}

\end{frame}

% =========================================================
\begin{frame}[fragile]{Completar la gráfica}

    \begin{pythoncodedark}
        plt.axhline(0)
        plt.axvline(0)

        plt.xlim(-6, 6)
        plt.ylim(-6, 6)

        plt.xlabel("X")
        plt.ylabel("Y")

        plt.axis("equal")
        plt.grid()
        plt.legend()

        plt.show()
    \end{pythoncodedark}

\end{frame}

% =========================================================
\begin{frame}{Completar la gráfica}

    \Idea{
        Una gráfica permite detectar rápidamente errores de signo,
        cuadrante o sentido de rotación.
    }

    \vspace{0.3cm}

    \Result{
        Si el resultado gráfico no coincide con la interpretación geométrica,
        es probable que exista un error en el signo del ángulo o en la matriz.
    }

\end{frame}

% =========================================================
\section{Laboratorio}

\SectionSlide
{Laboratorio de la sesión}
{Calcular, transformar y visualizar}

% =========================================================
\begin{frame}{Parte 1: desarrollo manual}

    Considere el vector:

    \[
    \mathbf{p}
    =
    \begin{bmatrix}
        3\\
        -2
    \end{bmatrix}
    \]

    Realice manualmente las siguientes rotaciones:

    \[
    \theta_1=30^\circ
    \]

    \[
    \theta_2=120^\circ
    \]

    \[
    \theta_3=-60^\circ
    \]

\end{frame}

% =========================================================
\begin{frame}{Parte 1: desarrollo manual}

    Para cada caso:

    \begin{enumerate}
        \item convierta el ángulo a radianes;
        \item construya $R(\theta)$;
        \item calcule $\mathbf{p}'=R(\theta)\mathbf{p}$;
        \item calcule la magnitud antes y después;
        \item determine el cuadrante del vector resultante.
    \end{enumerate}

\end{frame}

% =========================================================
\begin{frame}{Parte 2: composición de rotaciones}

    Para:

    \[
    \mathbf{p}
    =
    \begin{bmatrix}
        2\\
        1
    \end{bmatrix}
    \]

    aplique consecutivamente:

    \[
    \theta_1=25^\circ
    \]

    y:

    \[
    \theta_2=50^\circ
    \]

\end{frame}

% =========================================================
\begin{frame}{Parte 2: composición de rotaciones}

    Calcule manualmente:

    \begin{enumerate}
        \item $R(\theta_1)$;
        \item $R(\theta_2)$;
        \item $R(\theta_2)R(\theta_1)$;
        \item el vector resultante;
        \item una única matriz $R(\theta_1+\theta_2)$;
        \item compare ambos resultados.
    \end{enumerate}

\end{frame}

% =========================================================
\begin{frame}{Parte 3: sistema local y global}

    Un robot posee la pose:

    \[
    {}^{G}\mathbf{x}_R
    =
    \begin{bmatrix}
        4\\
        3\\
        60^\circ
    \end{bmatrix}
    \]

    Un sensor detecta un objeto en:

    \[
    {}^{R}\mathbf{p}
    =
    \begin{bmatrix}
        3\\
        1
    \end{bmatrix}
    \]

\end{frame}

% =========================================================
\begin{frame}{Parte 3: sistema local y global}

    Calcule manualmente:

    \begin{enumerate}
        \item la matriz $R(60^\circ)$;
        \item el punto local rotado;
        \item la posición global del objeto;
        \item la distancia entre el robot y el objeto.
    \end{enumerate}

    Utilice:

    \[
    {}^{G}\mathbf{p}
    =
    {}^{G}\mathbf{p}_R+
    R(\theta){}^{R}\mathbf{p}
    \]

\end{frame}

% =========================================================
\begin{frame}{Parte 4: transformación inversa}

    Considere:

    \[
    {}^{G}\mathbf{p}_R
    =
    \begin{bmatrix}
        -2\\
        4
    \end{bmatrix}
    \]

    \[
    \theta=135^\circ
    \]

    y un objeto ubicado globalmente en:

    \[
    {}^{G}\mathbf{p}
    =
    \begin{bmatrix}
        1\\
        7
    \end{bmatrix}
    \]

\end{frame}

% =========================================================
\begin{frame}{Parte 4: transformación inversa}

    Determine sus coordenadas respecto al robot utilizando:

    \[
    {}^{R}\mathbf{p}
    =
    R^T(\theta)
    \left(
    {}^{G}\mathbf{p}
    -
    {}^{G}\mathbf{p}_R
    \right)
    \]

    \vspace{0.3cm}

    Finalmente, interprete el resultado geométricamente.

\end{frame}

% =========================================================
\begin{frame}{Parte 5: validación computacional}

    Después de realizar los cálculos manuales:

    \begin{enumerate}
        \item implemente cada ejercicio con NumPy;
        \item compare los resultados manuales y computacionales;
        \item calcule el error absoluto entre ambos resultados;
    \end{enumerate}

\end{frame}

% =========================================================
\begin{frame}{Parte 5: validación computacional}

    Además:

    \begin{enumerate}
        \setcounter{enumi}{3}
        \item grafique los vectores antes y después de cada rotación;
        \item grafique el sistema global y el sistema local del robot;
        \item represente gráficamente los puntos transformados.
    \end{enumerate}

    \vspace{0.3cm}

    \Practice{
        El código debe implementar las expresiones matemáticas estudiadas.
        Python se utilizará para comprobar y visualizar los resultados,
        no para sustituir el desarrollo manual.
    }

\end{frame}

% =========================================================
\section{Cierre}

\SectionSlide
{Cierre de la sesión}
{De la geometría a los sistemas de referencia}

% =========================================================
\begin{frame}{¿Qué construimos hoy?}

    Partimos de un vector:

    \[
    \mathbf{p}
    =
    \begin{bmatrix}
        x\\
        y
    \end{bmatrix}
    \]

    y construimos matemáticamente:

    \[
    \boxed{
        R(\theta)
        =
        \begin{bmatrix}
            \cos\theta & -\sin\theta\\
            \sin\theta & \cos\theta
    \end{bmatrix}}
    \]

\end{frame}

% =========================================================
\begin{frame}{¿Qué construimos hoy?}

    Luego utilizamos esta matriz para:

    \begin{itemize}
        \item rotar vectores;
        \item componer rotaciones;
        \item invertir rotaciones;
        \item relacionar sistemas de referencia;
        \item transformar coordenadas locales a globales y viceversa.
    \end{itemize}

\end{frame}

% =========================================================
\begin{frame}{La ecuación más importante de la sesión}

    Para un robot ubicado en:

    \[
    {}^{G}\mathbf{p}_R
    =
    \begin{bmatrix}
        x_R\\
        y_R
    \end{bmatrix}
    \]

    con orientación $\theta$, un punto expresado localmente como:

    \[
    {}^{R}\mathbf{p}
    \]

    puede expresarse globalmente mediante:

    \[
    \boxed{
        {}^{G}\mathbf{p}
        =
        {}^{G}\mathbf{p}_R
        +
        R(\theta)
        {}^{R}\mathbf{p}
    }
    \]

\end{frame}

% =========================================================
\begin{frame}{La ecuación más importante de la sesión}

    \Idea{
        Esta relación será fundamental cuando el robot deba interpretar
        sensores, posición, odometría y objetivos de navegación utilizando
        diferentes sistemas de referencia.
    }

    \vspace{0.3cm}

    \Result{
        La ecuación combina orientación y posición para expresar un punto
        local dentro del sistema global.
    }

\end{frame}

% =========================================================
\begin{frame}{Preguntas de cierre}

    Antes de finalizar, responda:

    \begin{enumerate}
        \item ¿Por qué una rotación no cambia la magnitud de un vector?
        \item ¿Qué representa físicamente el ángulo $\theta$ de $R(\theta)$?
        \item ¿Por qué $R^{-1}(\theta)=R^T(\theta)$?
    \end{enumerate}

\end{frame}

% =========================================================
\begin{frame}{Preguntas de cierre}

    También responda:

    \begin{enumerate}
        \setcounter{enumi}{3}
        \item ¿Qué diferencia existe entre rotar y trasladar un punto?
        \item ¿Por qué un mismo punto puede tener coordenadas diferentes?
        \item ¿Qué diferencia existe entre ${}^{G}\mathbf{p}$ y
        ${}^{R}\mathbf{p}$?
    \end{enumerate}

    \vspace{0.3cm}

    \Result{
        Comprender los sistemas de referencia es indispensable para interpretar
        correctamente cualquier posición o movimiento de un robot.
    }

\end{frame}

% =========================================================
\begin{frame}{Próxima clase}

    En la siguiente sesión comenzaremos a utilizar:

    \[
    \boxed{\text{ROS 2 Humble}}
    \]

    para representar computacionalmente sistemas robóticos mediante:

    \begin{itemize}
        \item nodos;
        \item tópicos;
        \item mensajes;
        \item publicación y suscripción;
        \item movimiento de un robot mediante Turtlesim.
    \end{itemize}

\end{frame}

% =========================================================
\begin{frame}{Próxima clase}

    \Idea{
        La matemática seguirá siendo el modelo del sistema.
        ROS 2 será la infraestructura que nos permitirá implementar y comunicar
        sus componentes.
    }

\end{frame}

% =========================================================
\begin{frame}
    \centering
    \smartimage[0.8\linewidth]{Template/gracias.png}
\end{frame}


% =========================================================
\section{Solución del laboratorio}

\SectionSlide
{Solución del laboratorio}
{Desarrollo matemático y verificación de resultados}

% =========================================================
\begin{frame}{Parte 1: rotación de \texorpdfstring{$30^\circ$}{30 grados}}

    Para:

    \[
    \mathbf{p}
    =
    \begin{bmatrix}
        3\\
        -2
    \end{bmatrix},
    \qquad
    \theta_1=30^\circ
    \]

    \[
    R(30^\circ)
    =
    \begin{bmatrix}
        \frac{\sqrt{3}}{2} & -\frac{1}{2}\\
        \frac{1}{2} & \frac{\sqrt{3}}{2}
    \end{bmatrix}
    \]

    \[
    \mathbf{p}'
    =
    R(30^\circ)\mathbf{p}
    \]

    \[
    \mathbf{p}'
    =
    \begin{bmatrix}
        \frac{3\sqrt{3}}{2}+1\\
        \frac{3}{2}-\sqrt{3}
    \end{bmatrix}
    \]

    \[
    \boxed{
        \mathbf{p}'
        \approx
        \begin{bmatrix}
            3.598\\
            -0.232
    \end{bmatrix}}
    \]

\end{frame}

% =========================================================
\begin{frame}{Parte 1: interpretación para \texorpdfstring{$30^\circ$}{30 grados}}

    La magnitud original es:

    \[
    \|\mathbf{p}\|
    =
    \sqrt{3^2+(-2)^2}
    =
    \sqrt{13}
    \]

    La magnitud después de rotar es:

    \[
    \|\mathbf{p}'\|
    \approx
    \sqrt{3.598^2+(-0.232)^2}
    \approx
    \sqrt{13}
    \]

    \[
    \boxed{
        \|\mathbf{p}'\|=\|\mathbf{p}\|
    }
    \]

    \vspace{0.3cm}

    \Result{
        El vector resultante queda en el cuarto cuadrante, porque
        $x'>0$ y $y'<0$.
    }

\end{frame}

% =========================================================
\begin{frame}{Parte 1: rotación de \texorpdfstring{$120^\circ$}{120 grados}}

    Para:

    \[
    \theta_2=120^\circ
    \]

    \[
    R(120^\circ)
    =
    \begin{bmatrix}
        -\frac{1}{2} & -\frac{\sqrt{3}}{2}\\
        \frac{\sqrt{3}}{2} & -\frac{1}{2}
    \end{bmatrix}
    \]

    \[
    \mathbf{p}'
    =
    R(120^\circ)
    \begin{bmatrix}
        3\\
        -2
    \end{bmatrix}
    \]

    \[
    \mathbf{p}'
    =
    \begin{bmatrix}
        -\frac{3}{2}+\sqrt{3}\\
        1+\frac{3\sqrt{3}}{2}
    \end{bmatrix}
    \]

    \[
    \boxed{
        \mathbf{p}'
        \approx
        \begin{bmatrix}
            0.232\\
            3.598
    \end{bmatrix}}
    \]

\end{frame}

% =========================================================
\begin{frame}{Parte 1: interpretación para \texorpdfstring{$120^\circ$}{120 grados}}

    La magnitud se conserva:

    \[
    \|\mathbf{p}\|=\sqrt{13}
    \]

    \[
    \|\mathbf{p}'\|
    \approx
    \sqrt{0.232^2+3.598^2}
    \approx
    \sqrt{13}
    \]

    \vspace{0.3cm}

    \Result{
        El vector resultante queda en el primer cuadrante, porque
        $x'>0$ y $y'>0$.
    }

\end{frame}

% =========================================================
\begin{frame}{Parte 1: rotación de \texorpdfstring{$-60^\circ$}{-60 grados}}

    Para:

    \[
    \theta_3=-60^\circ
    \]

    \[
    R(-60^\circ)
    =
    \begin{bmatrix}
        \frac{1}{2} & \frac{\sqrt{3}}{2}\\
        -\frac{\sqrt{3}}{2} & \frac{1}{2}
    \end{bmatrix}
    \]

    \[
    \mathbf{p}'
    =
    R(-60^\circ)
    \begin{bmatrix}
        3\\
        -2
    \end{bmatrix}
    \]

    \[
    \mathbf{p}'
    =
    \begin{bmatrix}
        \frac{3}{2}-\sqrt{3}\\
        -\frac{3\sqrt{3}}{2}-1
    \end{bmatrix}
    \]

    \[
    \boxed{
        \mathbf{p}'
        \approx
        \begin{bmatrix}
            -0.232\\
            -3.598
    \end{bmatrix}}
    \]

\end{frame}

% =========================================================
\begin{frame}{Parte 1: interpretación para \texorpdfstring{$-60^\circ$}{-60 grados}}

    La magnitud también se conserva:

    \[
    \|\mathbf{p}\|=\sqrt{13}
    \]

    \[
    \|\mathbf{p}'\|
    \approx
    \sqrt{(-0.232)^2+(-3.598)^2}
    \approx
    \sqrt{13}
    \]

    \vspace{0.3cm}

    \Result{
        El vector resultante queda en el tercer cuadrante, porque
        $x'<0$ y $y'<0$.
    }

\end{frame}

% =========================================================
\begin{frame}{Parte 2: composición de rotaciones}

    Para:

    \[
    \mathbf{p}
    =
    \begin{bmatrix}
        2\\
        1
    \end{bmatrix},
    \qquad
    \theta_1=25^\circ,
    \qquad
    \theta_2=50^\circ
    \]

    La rotación equivalente es:

    \[
    \theta_T=\theta_1+\theta_2
    \]

    \[
    \theta_T=25^\circ+50^\circ
    \]

    \[
    \boxed{\theta_T=75^\circ}
    \]

    Por tanto:

    \[
    R(50^\circ)R(25^\circ)=R(75^\circ)
    \]

\end{frame}

% =========================================================
\begin{frame}{Parte 2: matrices de rotación}

    Aproximadamente:

    \[
    R(25^\circ)
    \approx
    \begin{bmatrix}
        0.9063 & -0.4226\\
        0.4226 & 0.9063
    \end{bmatrix}
    \]

    \[
    R(50^\circ)
    \approx
    \begin{bmatrix}
        0.6428 & -0.7660\\
        0.7660 & 0.6428
    \end{bmatrix}
    \]

    Al multiplicar:

    \[
    R(50^\circ)R(25^\circ)
    \approx
    \begin{bmatrix}
        0.2588 & -0.9659\\
        0.9659 & 0.2588
    \end{bmatrix}
    \]

\end{frame}

% =========================================================
\begin{frame}{Parte 2: comparación con \texorpdfstring{$R(75^\circ)$}{R 75}}

    La matriz directa para $75^\circ$ es:

    \[
    R(75^\circ)
    \approx
    \begin{bmatrix}
        0.2588 & -0.9659\\
        0.9659 & 0.2588
    \end{bmatrix}
    \]

    Por tanto:

    \[
    \boxed{
        R(50^\circ)R(25^\circ)=R(75^\circ)
    }
    \]

    Aplicando sobre el vector:

    \[
    \mathbf{p}'
    =
    R(75^\circ)
    \begin{bmatrix}
        2\\
        1
    \end{bmatrix}
    \]

    \[
    \boxed{
        \mathbf{p}'
        \approx
        \begin{bmatrix}
            -0.448\\
            2.191
    \end{bmatrix}}
    \]

\end{frame}

% =========================================================
\begin{frame}{Parte 3: sistema local y global}

    El robot posee la pose:

    \[
    {}^{G}\mathbf{x}_R
    =
    \begin{bmatrix}
        4\\
        3\\
        60^\circ
    \end{bmatrix}
    \]

    Por tanto:

    \[
    {}^{G}\mathbf{p}_R
    =
    \begin{bmatrix}
        4\\
        3
    \end{bmatrix}
    \]

    El sensor detecta:

    \[
    {}^{R}\mathbf{p}
    =
    \begin{bmatrix}
        3\\
        1
    \end{bmatrix}
    \]

    La matriz es:

    \[
    R(60^\circ)
    =
    \begin{bmatrix}
        \frac{1}{2} & -\frac{\sqrt{3}}{2}\\
        \frac{\sqrt{3}}{2} & \frac{1}{2}
    \end{bmatrix}
    \]

\end{frame}

% =========================================================
\begin{frame}{Parte 3: punto local rotado}

    Aplicamos la rotación:

    \[
    R(60^\circ)
    {}^{R}\mathbf{p}
    =
    \begin{bmatrix}
        \frac{1}{2} & -\frac{\sqrt{3}}{2}\\
        \frac{\sqrt{3}}{2} & \frac{1}{2}
    \end{bmatrix}
    \begin{bmatrix}
        3\\
        1
    \end{bmatrix}
    \]

    \[
    =
    \begin{bmatrix}
        \frac{3}{2}-\frac{\sqrt{3}}{2}\\
        \frac{3\sqrt{3}}{2}+\frac{1}{2}
    \end{bmatrix}
    \]

    \[
    \boxed{
        R(60^\circ)
        {}^{R}\mathbf{p}
        \approx
        \begin{bmatrix}
            0.634\\
            3.098
    \end{bmatrix}}
    \]

\end{frame}

% =========================================================
\begin{frame}{Parte 3: posición global del objeto}

    La transformación local a global es:

    \[
    {}^{G}\mathbf{p}
    =
    {}^{G}\mathbf{p}_R
    +
    R(\theta)
    {}^{R}\mathbf{p}
    \]

    Sustituyendo:

    \[
    {}^{G}\mathbf{p}
    =
    \begin{bmatrix}
        4\\
        3
    \end{bmatrix}
    +
    \begin{bmatrix}
        0.634\\
        3.098
    \end{bmatrix}
    \]

    \[
    \boxed{
        {}^{G}\mathbf{p}
        \approx
        \begin{bmatrix}
            4.634\\
            6.098
    \end{bmatrix}}
    \]

\end{frame}

% =========================================================
\begin{frame}{Parte 3: distancia entre robot y objeto}

    La distancia entre el robot y el objeto es la magnitud del vector local:

    \[
    {}^{R}\mathbf{p}
    =
    \begin{bmatrix}
        3\\
        1
    \end{bmatrix}
    \]

    \[
    d
    =
    \sqrt{3^2+1^2}
    \]

    \[
    \boxed{
        d=\sqrt{10}
    }
    \]

    \[
    \boxed{
        d\approx3.162
    }
    \]

    \vspace{0.3cm}

    \Result{
        La rotación cambia las coordenadas del punto respecto al sistema global,
        pero no cambia la distancia entre el robot y el objeto.
    }

\end{frame}

% =========================================================
\begin{frame}{Parte 4: transformación inversa}

    Se tiene:

    \[
    {}^{G}\mathbf{p}_R
    =
    \begin{bmatrix}
        -2\\
        4
    \end{bmatrix}
    \]

    \[
    \theta=135^\circ
    \]

    \[
    {}^{G}\mathbf{p}
    =
    \begin{bmatrix}
        1\\
        7
    \end{bmatrix}
    \]

    Primero restamos:

    \[
    {}^{G}\mathbf{p}
    -
    {}^{G}\mathbf{p}_R
    =
    \begin{bmatrix}
        1\\
        7
    \end{bmatrix}
    -
    \begin{bmatrix}
        -2\\
        4
    \end{bmatrix}
    =
    \begin{bmatrix}
        3\\
        3
    \end{bmatrix}
    \]

\end{frame}

% =========================================================
\begin{frame}{Parte 4: matriz inversa}

    La transformación inversa utiliza:

    \[
    {}^{R}\mathbf{p}
    =
    R^T(135^\circ)
    \left(
    {}^{G}\mathbf{p}
    -
    {}^{G}\mathbf{p}_R
    \right)
    \]

    Como:

    \[
    R^T(135^\circ)=R(-135^\circ)
    \]

    \[
    R(-135^\circ)
    =
    \begin{bmatrix}
        -\frac{\sqrt{2}}{2} & \frac{\sqrt{2}}{2}\\
        -\frac{\sqrt{2}}{2} & -\frac{\sqrt{2}}{2}
    \end{bmatrix}
    \]

\end{frame}

% =========================================================
\begin{frame}{Parte 4: coordenadas respecto al robot}

    Entonces:

    \[
    {}^{R}\mathbf{p}
    =
    \begin{bmatrix}
        -\frac{\sqrt{2}}{2} & \frac{\sqrt{2}}{2}\\
        -\frac{\sqrt{2}}{2} & -\frac{\sqrt{2}}{2}
    \end{bmatrix}
    \begin{bmatrix}
        3\\
        3
    \end{bmatrix}
    \]

    \[
    {}^{R}\mathbf{p}
    =
    \begin{bmatrix}
        0\\
        -3\sqrt{2}
    \end{bmatrix}
    \]

    \[
    \boxed{
        {}^{R}\mathbf{p}
        \approx
        \begin{bmatrix}
            0\\
            -4.243
    \end{bmatrix}}
    \]

\end{frame}

% =========================================================
\begin{frame}{Parte 4: interpretación geométrica}

    El resultado fue:

    \[
    {}^{R}\mathbf{p}
    \approx
    \begin{bmatrix}
        0\\
        -4.243
    \end{bmatrix}
    \]

    Esto significa que, visto desde el sistema del robot:

    \begin{itemize}
        \item el objeto no está desplazado hacia adelante ni hacia atrás en el eje local $x_R$;
        \item el objeto se encuentra sobre el eje local negativo $y_R$;
        \item la distancia al objeto es aproximadamente $4.243$ unidades.
    \end{itemize}

    \Result{
        El punto global queda expresado ahora desde la perspectiva del robot.
    }

\end{frame}
\end{document}